\documentclass[11pt,a4paper]{article}

\usepackage{amsmath,amssymb,amsthm}
\usepackage{bm}
\usepackage{mathtools}
\usepackage{physics}
\usepackage[colorlinks=true,linkcolor=blue,urlcolor=blue,citecolor=blue]{hyperref}

\DeclareMathOperator{\E}{\mathbb{E}}
\DeclareMathOperator{\svd}{SVD}

\title{LoRC: Low-rank Quantization Correction}
\date{}

\begin{document}
\maketitle

\section{Introduction}

We consider the problem of compressing a pretrained language model
$f_\theta$ while preserving its performance on a target domain
$\mathcal{D}_{\text{target}}$ at the expense of out-of-domain
generality.  Our approach -- \textbf{Lo}w-rank \textbf{Q}uantization
\textbf{C}orrection -- identifies low-dimensional subspaces where
activation statistics differ between two domains and uses them to
correct quantization error.

\subsection{Covariance estimation}

For a linear layer with weight $W \in \mathbb{R}^{d_{\text{out}} \times
d_{\text{in}}}$, let $x \in \mathbb{R}^{d_{\text{in}}}$ and $y \in
\mathbb{R}^{d_{\text{out}}}$ be its pre- and post-activation
respectively.  For each domain $\mathcal{D}$, we collect the
second-moment statistics

\begin{align}
  C_{\text{pre}}^{\mathcal{D}} &= \E[xx^{\mathsf{T}}] -
  \E[x]\E[x]^{\mathsf{T}} \in \mathbb{R}^{d_{\text{in}} \times
  d_{\text{in}}}, \\
  C_{\text{post}}^{\mathcal{D}} &= \E[yy^{\mathsf{T}}] -
  \E[y]\E[y]^{\mathsf{T}} \in \mathbb{R}^{d_{\text{out}} \times
  d_{\text{out}}},
\end{align}

estimated by accumulating $\sum xx^{\mathsf{T}}$ and $\sum x$ over a
stream of domain-specific prompts.

\subsection{Domain subspace extraction}

Given two domains $\mathcal{D}_1$ and $\mathcal{D}_2$ (e.g., Lean4
proofs and Wikipedia text), form the difference

\begin{align}
  \Delta C &= \alpha\,C^{\mathcal{D}_1} - \beta\,C^{\mathcal{D}_2},
  \qquad \alpha,\beta \ge 0,
\end{align}

symmetrize as $(\Delta C + \Delta C^{\mathsf{T}})/2$, and compute the
eigendecomposition

\begin{align}
  \Delta C &= Q \Lambda Q^{\mathsf{T}}.
\end{align}

The eigenvectors with the largest eigenvalues correspond to
directions where $\mathcal{D}_1$ activates more than $\mathcal{D}_2$
(the $\mathcal{D}_1$-dominant subspace); those with the most negative
eigenvalues correspond to the $\mathcal{D}_2$-dominant subspace.  We
keep $K \ll d$ eigenvectors in each direction, forming the
column-orthonormal matrices $V_1, V_2 \in \mathbb{R}^{d \times K}$.

\subsection{Correction factors}

Quantize the weight matrix to NF4 and compute the quantization error

\begin{align}
  E = W - Q^{-1}\bigl(Q(W)\bigr).
\end{align}

Project $E$ onto the $\mathcal{D}_1$ subspace (the $\mathcal{D}_2$ case
is symmetric).  For a weight applied on the input side (e.g.\
up-projection), the projection is

\begin{align}
  E_{\text{proj}} = E\, (V_1 V_1^{\mathsf{T}}),
\end{align}

and for a weight applied on the output side (e.g.\ down-projection)

\begin{align}
  E_{\text{proj}} = (V_1 V_1^{\mathsf{T}})\, E.
\end{align}

Factorize $E_{\text{proj}} \in \mathbb{R}^{d_{\text{out}} \times
d_{\text{in}}}$ via the singular value decomposition

\begin{align}
  U \Sigma V^{\mathsf{T}} &= \svd(E_{\text{proj}}), \\
  U_k \Sigma_k V_k^{\mathsf{T}} &\approx E_{\text{proj}},
\end{align}

truncated to rank $k = \min(K, d_{\text{out}}, d_{\text{in}})$.
Define the correction factors

\begin{align}
  V_{\text{act}} &= V_k \Sigma_k^{1/2} \in \mathbb{R}^{d_{\text{in}} \times k}, \\
  U_{\text{write}} &= U_k \Sigma_k^{1/2} \in \mathbb{R}^{d_{\text{out}} \times k},
\end{align}

both stored in bfloat16.

\subsection{Hybrid forward pass}

At inference, the base weight remains in NF4.  The forward pass becomes

\begin{align}
  z = x W_{\text{base}}^{\mathsf{T}} + (x V_{\text{act}}) U_{\text{write}}^{\mathsf{T}},
\end{align}

where the first term is computed via on-the-fly dequantization of the
4-bit weights and the second term supplies the low-rank BF16
correction for the active domain.

\subsection{Disjunction score}

To measure whether the identified subspaces are genuinely
domain-specific, we project them out of the original weight matrix.
For weights applied on the input side (e.g.\ up-projection):

\begin{align}
  W' = W - W\,(V V^{\mathsf{T}}),
\end{align}

and for weights applied on the output side (e.g.\ down-projection):

\begin{align}
  W' = W - (V V^{\mathsf{T}})\, W.
\end{align}

We then measure the perplexity change $\Delta\text{PPL}$ on each domain.
The disjunction score

\begin{align}
  D = \frac{\max(\Delta\text{PPL}_{\text{own}},\ \varepsilon)}
           {\max(\Delta\text{PPL}_{\text{other}},\ \varepsilon)}
\end{align}

quantifies domain specificity: $D \gg 1$ indicates genuine separation,
while $D \approx 1$ indicates polysemantic (entangled) directions.

\subsection{Possible improvements}

The current implementation incurs $\mathcal{O}(d^2)$ memory for
covariance storage and $\mathcal{O}(d^3)$ for the
eigendecompositions, which limits scalability to larger models.
Several conceptually motivated optimizations exist.

\paragraph{Streaming PCA.}
Instead of accumulating the full outer-product matrix $\sum xx^{\mathsf{T}}$,
maintain a subspace estimate $V \in \mathbb{R}^{d \times K}$ via
online updates (Oja's rule):
\begin{align}
  V \leftarrow V + \eta\,\bigl( x x^{\mathsf{T}} V \bigr), \qquad
  V \leftarrow \text{QR}(V),
\end{align}
where $x$ is a sampled activation.  Storage drops from
$\mathcal{O}(d^2)$ to $\mathcal{O}(dK)$, and the two domains can be
processed concurrently using domain labels.

\paragraph{Randomized eigendecomposition.}
Only the $K$ extremal eigenvalues of $\Delta C$ are needed.  Replace
the full $\text{eigh}(\Delta C)$ with a Krylov subspace method:
\begin{align}
  \Omega &\sim \mathcal{N}(0,1)^{d \times (K + p)}, &
  Y &= \Delta C \,\Omega, &
  Q &= \text{QR}(Y), \\
  T &= Q^{\mathsf{T}} \Delta C \, Q \in \mathbb{R}^{(K+p) \times (K+p)}, &
  \tilde\Lambda, \tilde U &= \text{eigh}(T), &
  V_{\text{ext}} &= Q \tilde U,
\end{align}
with oversampling $p \approx 5$ for numerical stability.  The cost is
$\mathcal{O}(K d^2)$ rather than $\mathcal{O}(d^3)$.

\paragraph{Avoiding the explicit projection matrix.}
The projected error $E_{\text{proj}} = E (V V^{\mathsf{T}})$ is never
needed explicitly.  Factorize $A = E V \in \mathbb{R}^{d_{\text{out}}
\times K}$ directly:
\begin{align}
  U_A \Sigma_A R_A^{\mathsf{T}} &= \svd(A), \\
  V_{\text{act}} &= (V R_A) \Sigma_A^{1/2}, \qquad
  U_{\text{write}} = U_A \Sigma_A^{1/2}.
\end{align}
This produces identical correction factors without constructing the
$d_{\text{in}} \times d_{\text{in}}$ projection matrix and without a
full-size SVD.  Total cost reduces from $\mathcal{O}(d^3 + d^2 K)$ to
$\mathcal{O}(d^2 K)$.

\paragraph{Analytic causal masking.}
The singular values $\Sigma_A$ of $A = E V$ are the
per-component contributions to the projected quantization error.
Gradient-based mask training can be replaced by a simple threshold:
\begin{align}
  \text{keep}_i = \mathbf{1}\bigl( \sigma_i > \tau \bigr),
  \qquad
  \tau = \text{quantile}\bigl(\Sigma,\; 1 - \text{keep\_pct}\bigr),
\end{align}
eliminating the data-dependent optimization loop entirely.

\paragraph{Extension to more than two domains.}
For $D > 2$ domains, replace the pairwise difference $\Delta C$ with
multi-class Fisher discriminant analysis:
\begin{align}
  W^* = \arg\max_{W^{\mathsf{T}} W = I}
  \frac{\tr\bigl( W^{\mathsf{T}} C_{\text{between}} W \bigr)}
       {\tr\bigl( W^{\mathsf{T}} C_{\text{within}} W \bigr)},
\end{align}
where $C_{\text{between}}$ is the between-class scatter matrix and
$C_{\text{within}}$ is the pooled within-class covariance.  The
optimal $W$ is given by the generalized eigenvectors of
$C_{\text{between}}\, v = \lambda C_{\text{within}}\, v$, yielding a
single subspace that maximally separates all $D$ domains.

\paragraph{Cross-layer subspace coupling.}
Per-module subspaces are estimated independently.  Constraining them
to share a global basis $V_{\text{gl}} \in \mathbb{R}^{d \times K}$
with per-module adapters $M_{\ell}$ such that
$V_{\ell} = V_{\text{gl}} M_{\ell}$ reduces overfitting and total
parameter count, especially in deeper models where domain-specific
patterns recur across layers.

\end{document}
